% Figure: the isobaric heat capacity c_p of a supercritical fluid against
% temperature at three supercritical pressures. Near the critical pressure
% the peak is tall and sharp; at higher pressure it broadens and shifts to a
% higher temperature. The locus of the peaks is the Widom line, the soft
% continuation of the saturation curve above the critical point. Schematic.
\begin{figure}[htbp]
\centering
\begin{tikzpicture}
\begin{axis}[
  width=12cm, height=7.5cm,
  xlabel={temperature $T$ (\si{\degreeCelsius})},
  ylabel={isobaric heat capacity $c_p$ (arb.\ units)},
  xmin=25, xmax=75, ymin=0, ymax=10,
  axis lines=left,
  tick label style={font=\scriptsize},
  label style={font=\small},
  legend pos=north east,
  legend style={font=\scriptsize},
  clip=true,
]
% p = 8.0 MPa (just above p_c): tall sharp peak near 34 C
\addplot[engred, very thick, smooth, samples=120, domain=25:75]
  {1.2 + 8.2*exp(-((x-34)^2)/(2*2.2^2))};
\addlegendentry{$p=8.0$ MPa (near $p_c$)}
% p = 9.5 MPa: lower broader peak near 40 C
\addplot[engorange, very thick, smooth, samples=120, domain=25:75]
  {1.2 + 4.6*exp(-((x-40)^2)/(2*4.0^2))};
\addlegendentry{$p=9.5$ MPa}
% p = 12.0 MPa: broad low peak near 48 C
\addplot[engblue, very thick, smooth, samples=120, domain=25:75]
  {1.2 + 2.6*exp(-((x-48)^2)/(2*6.5^2))};
\addlegendentry{$p=12.0$ MPa}
% Widom line: connect the three peak locations
\addplot[enggray, dashed, thick] coordinates {(34,9.4) (40,5.8) (48,3.8)};
\node[anchor=west, font=\scriptsize, enggray] at (axis cs:48.5,3.9)
  {Widom line (locus of $c_p$ maxima)};
\end{axis}
\end{tikzpicture}
\caption[The supercritical heat-capacity peak and the Widom line]{The
isobaric heat capacity $c_p$ of carbon dioxide against temperature at three
supercritical pressures. Just above the critical pressure the peak is tall
and sharp; raising the pressure lowers and broadens it and shifts it to a
higher temperature. The peak carries no latent heat, because there is no
phase transition above the critical point, but a heat exchanger cooling its
fluid across the peak passes through a region of large effective heat
capacity that acts like a smeared-out latent heat. The dashed locus joining
the peaks is the Widom line, the soft continuation of the saturation curve
into the single-phase region: below it the supercritical fluid is dense and
liquid-like, above it dilute and gas-like. A transcritical CO\textsubscript{2}
gas cooler rejects its heat by carrying the high-side fluid across this peak,
which is why its outlet temperature matters so much to the cycle.}
\label{fig:vol04ch02:widom-cp-peak}
\end{figure}
